% Standalone LaTeX edition of neutron-diffusion-paper.md.
% Compile with Tectonic or run pdflatex twice for references.
\documentclass[11pt,letterpaper]{article}
\usepackage[margin=0.85in]{geometry}
\usepackage{amsmath,amssymb}
\usepackage{booktabs,tabularx}
\usepackage{xcolor}
\usepackage{microtype}
\usepackage{hyperref}

\definecolor{linkblue}{HTML}{18577B}
\hypersetup{
  colorlinks=true,
  linkcolor=linkblue,
  citecolor=linkblue,
  urlcolor=linkblue,
  pdftitle={Two-group neutron diffusion in the Candu Refuelling Simulator},
  pdfauthor={CRS technical paper}
}
\setlength{\emergencystretch}{2em}
\setcounter{secnumdepth}{2}

\title{\textbf{Two-group neutron diffusion}\\
       \textbf{in the Candu Refuelling Simulator}\\[0.35em]
       \large Numerical formulation, verification, and browser latency}
\author{CRS technical paper}
\date{9 November 2026 \enspace\textperiodcentered\enspace Runtime revision \texttt{4341ae8}}

\begin{document}
\maketitle

\begin{abstract}
The Candu Refuelling Simulator (CRS) computes the steady neutron flux and effective
multiplication factor of a stylized CANDU-6 core in a browser. The model has 4,560
finite volumes, two energy groups, and a sparse leakage stencil. The solver applies
power iteration with two inexact Jacobi solves per outer iteration and normalizes the
resulting flux to prescribed thermal power. Reciprocal face currents conserve neutrons
across interior faces, and strict row diagonal dominance gives a convergence argument
for the fixed-source group solves. An analytic reflective-cell solution and spatial
regressions verify aspects of the discrete model. For gameplay, one frozen-flux
integration covers each simulated half-day, with at most one additional equilibrium
event for the entire refuelling batch. Warm starts, reused candidate inputs, and
ahead-of-time WebAssembly compilation further reduce command latency. In a desktop
Chromium run, all 27 measured AOT commands finished within five seconds. Median times
were 1.91 s, 3.82 s, and 3.94 s for zero, one, and two fuel orders respectively.
\end{abstract}

\noindent\textbf{Keywords.} neutron diffusion, finite volume, source iteration, Jacobi, WebAssembly, CANDU,
numerical verification

\section{Problem and scope}\label{problem-and-scope}

After a player replaces fuel, CRS must recompute where fission power is produced and how
the liquid-zone regulator responds. A local fuel change affects neighboring channels
through leakage, so each equilibrium event solves the whole configured core. The browser
requests steps from the same authoritative C\# simulation used by the native tests. Core
owns neutronics and inventory, Game owns the transaction and operating-limit decisions,
and the browser renders accepted observations.
\href{https://github.com/hkkho/CRS/blob/4341ae8881f5827988c3d8b1859aa9e8f13cd8a6/docs/IMPLEMENTATION_GUIDE.md}{Architecture
and cadence} ·
\href{https://github.com/hkkho/CRS/blob/4341ae8881f5827988c3d8b1859aa9e8f13cd8a6/web/candu-playtest/src/wasmWorker.ts}{Browser
worker}

The default topology follows the 380-channel, 12-bundle geometry of the CANDU-6
reference design. CRS homogenizes each bundle position into one volume with two scalar
flux unknowns, giving 9,120 flux unknowns and one effective multiplication factor. The
active v11 pack is a calibrated gameplay surrogate with authored coefficients, effective
boundaries and localized absorption. This paper assesses the balance equations of that
surrogate and the cost of solving them. Plant-level predictive accuracy would require
independent reference data and model-error studies. \cite{ref-1} ·
\href{https://github.com/hkkho/CRS/blob/4341ae8881f5827988c3d8b1859aa9e8f13cd8a6/data/packs/candu6-two-group-diffusion-pack-v1.json}{Active
data pack} ·
\href{https://github.com/hkkho/CRS/blob/4341ae8881f5827988c3d8b1859aa9e8f13cd8a6/docs/physics/calibration-provenance.md}{Calibration
provenance}

\section{Physical and spatial formulation}\label{physical-and-spatial-formulation}

\subsection{Steady two-group balance}\label{steady-two-group-balance}

The diffusion approximation closes the angular transport balance with
Fick\textquotesingle s law, \(\mathbf J_g=-D_g\nabla\phi_g\). Here \(\phi_g\) is
group-integrated scalar flux in neutrons m\(^{-2}\) s\(^{-1}\), \(D_g\) is an effective
diffusion coefficient in metres, and \(\mathbf J_g\) is the net current density. Index 1
denotes fast neutrons and index 2 thermal neutrons. With the material state frozen,
downscatter transfers neutrons from group 1 to group 2, and the model neglects
upscatter. The steady production, removal and leakage equations are \cite{ref-2} ·
\cite{ref-3}

\begin{equation}
\label{eq:group-balance}
\begin{aligned}
-\nabla\cdot(D_1\nabla\phi_1)+(\Sigma_{a1}+\Sigma_{s12})\phi_1
 &= \chi_1 F/k,\\
-\nabla\cdot(D_2\nabla\phi_2)+\Sigma_{a2}\phi_2
 &= \Sigma_{s12}\phi_1+\chi_2 F/k,\\
F &= \nu\Sigma_{f1}\phi_1+\nu\Sigma_{f2}\phi_2,
\qquad \chi_1+\chi_2=1.
\end{aligned}
\end{equation}

Macroscopic cross sections have units m\(^{-1}\), so every term in the balance is a
volumetric rate. \(F\) counts fission neutrons produced per volume per second. The
default fuel rows put all fission births in the fast group. The eigenvalue \(k\) scales
production to achieve steady balance. A configuration at \(k=1\) is critical in the
model, and the associated dimensionless reactivity is \(\rho=(k-1)/k\). Because the
eigenproblem is homogeneous in flux, the amplitude is fixed separately using the target
thermal power \(P_\star\). The solver computes the heat in cell \(i\) from its fission
rate and energy per fission \(E_{f,i}\).
\href{https://github.com/hkkho/CRS/blob/4341ae8881f5827988c3d8b1859aa9e8f13cd8a6/src/ReactorSim.Core/Domain/FullCoreDiffusionModelContracts.cs\#L1480}{Pack
binding} ·
\href{https://github.com/hkkho/CRS/blob/4341ae8881f5827988c3d8b1859aa9e8f13cd8a6/src/ReactorSim.Core/Domain/SpatialEigenIterationContracts.cs\#L831}{Eigen
iteration and power}

\begin{equation}
\label{eq:thermal-power}
P_i=V_iE_{f,i}\big(\Sigma_{f1,i}\phi_{1,i}+\Sigma_{f2,i}\phi_{2,i}\big),
\qquad P(\phi)=\sum_i P_i=P_\star .
\end{equation}

The default \(P_\star\) is 2,064 MW thermal. The heating coefficients \(\Sigma_f\) and
the neutron-production coefficients \(\nu\Sigma_f\) enter different equations. A
configuration may have \(k\ne 1\) at the prescribed power because normalization sets the
amplitude of a steady shape. The model uses that shape to calculate subsequent burnup
and poison exposure.
\href{https://github.com/hkkho/CRS/blob/4341ae8881f5827988c3d8b1859aa9e8f13cd8a6/docs/physics/active-two-group-solver.md}{Active
solver and coefficients} ·
\href{https://github.com/hkkho/CRS/blob/4341ae8881f5827988c3d8b1859aa9e8f13cd8a6/docs/physics/channel-power-reference.md}{Reference
power}

\subsection{Conservative finite volumes}\label{conservative-finite-volumes}

Integrate \eqref{eq:group-balance} over cell \(i\), apply the divergence theorem and
divide by its volume \(V_i\). Approximating the integrated face current by a conductance
times a two-point flux difference gives the group operator \cite{ref-3} ·
\href{https://github.com/hkkho/CRS/blob/4341ae8881f5827988c3d8b1859aa9e8f13cd8a6/src/ReactorSim.Core/Domain/SpatialStencilContracts.cs\#L116}{Stencil
assembly} ·
\href{https://github.com/hkkho/CRS/blob/4341ae8881f5827988c3d8b1859aa9e8f13cd8a6/src/ReactorSim.Core/Domain/SpatialOperatorContracts.cs\#L1213}{Operator
evaluation}

\begin{equation}
\label{eq:finite-volume}
\begin{aligned}
(A_g\phi_g)_i &= R_{g,i}\phi_{g,i}+\frac{1}{V_i}\left[
\sum_{j\sim i} C_{g,ij}(\phi_{g,i}-\phi_{g,j})
+\sum_{b\in\partial i} C_{g,ib}\phi_{g,i}
\right],\\
R_{1,i} &= \Sigma_{a1,i}+\Sigma_{s12,i},
\qquad R_{2,i}=\Sigma_{a2,i}.
\end{aligned}
\end{equation}

For an interior face, the two nodes share the same \(C_{g,ij}=C_{g,ji}\). Multiplying
each row by its cell volume before summing cancels these two currents exactly in the
algebra. Summing the two groups then cancels downscatter as an internal transfer. The
remaining whole-core balance is fission production divided by \(k\) equal to absorption
plus outward boundary leakage. Conductances are in m\(^2\), consistent with a face area
times diffusion coefficient divided by distance.
\href{https://github.com/hkkho/CRS/blob/4341ae8881f5827988c3d8b1859aa9e8f13cd8a6/src/ReactorSim.Core/Domain/SpatialOperatorContracts.cs\#L1011}{Conductance
binding} ·
\href{https://github.com/hkkho/CRS/blob/4341ae8881f5827988c3d8b1859aa9e8f13cd8a6/docs/physics/active-two-group-solver.md}{Boundary
and conductance values}

The pack uses effective diffusion coefficients \(D_1=0.015\) m and \(D_2=0.0075\) m,
with \(V=0.04044276185625\) m\(^3\) per node. Interior conductance is of the form
\(D_g A_f/d\) for face area \(A_f\) and node-centre separation \(d\). At a radial vacuum
face the distance is from the node centre to the 5.55 m extrapolation circle along the
outward coordinate ray. At an axial vacuum face the distance is half the 0.4953 m cell
length plus 0.22 m. The extrapolated flux is zero in both cases. A reflective face has
zero conductance. These are effective leakage closures on the channel lattice.
\href{https://github.com/hkkho/CRS/blob/4341ae8881f5827988c3d8b1859aa9e8f13cd8a6/data/calibration/diffusion-v10/README.md}{Geometry
calibration} ·
\href{https://github.com/hkkho/CRS/blob/4341ae8881f5827988c3d8b1859aa9e8f13cd8a6/src/ReactorSim.Core/Domain/CircularRadialBoundary.cs}{Radial
boundary implementation}

Fuel coefficients are bound by bundle burnup from the pack\textquotesingle s
piecewise-linear tables. Fixed adjusters and liquid-zone water change local absorption.
Live xenon replaces the xenon already included in the fuel table by adding
\((N_{\rm Xe}-N_{\rm Xe,ref})\sigma_{\rm capture}w\) to thermal absorption, where \(w\)
is the pack\textquotesingle s separate worth scale. The spatial solve holds these
coefficients fixed for each candidate.
\href{https://github.com/hkkho/CRS/blob/4341ae8881f5827988c3d8b1859aa9e8f13cd8a6/src/ReactorSim.Core/Domain/BurnupCoefficientContracts.cs}{Burnup
lookup} ·
\href{https://github.com/hkkho/CRS/blob/4341ae8881f5827988c3d8b1859aa9e8f13cd8a6/docs/physics/xenon-calibration.md}{Xenon
calibration}

\section{Eigenvalue algorithm and numerical
correctness}\label{eigenvalue-algorithm-and-numerical-correctness}

\subsection{Source iteration}\label{source-iteration}

Let \(A\) be block lower triangular with \(A_1\) and \(A_2\) on the diagonal and
\(-S_{12}\) below it, where \(S_{12}\) is the diagonal downscatter matrix. Let
\(Q\phi=F\) compute the cellwise fission neutron source, and let \(X\) distribute it
into the two groups according to \(\chi\). Equation \eqref{eq:group-balance} after
discretization is \(A\phi=XQ\phi/k\). One outer step forms \(F\) from the current flux,
approximately solves \(A_1\) with source \(\chi_1F/k\), then approximately solves
\(A_2\) with source \(S_{12}\widetilde\phi_1+
\chi_2F/k\) using the trial fast flux \(\widetilde\phi_1\). With
\(T(\phi)=\sum_i V_i(Q\phi)_i\), the implementation updates \(k\) before normalizing
both groups by the same factor.
\href{https://github.com/hkkho/CRS/blob/4341ae8881f5827988c3d8b1859aa9e8f13cd8a6/src/ReactorSim.Core/Domain/SpatialEigenIterationContracts.cs\#L388}{Source
iteration implementation}

\begin{equation}
\label{eq:iteration}
k_{\rm next}=k\,\frac{T(\widetilde\phi)}{T(\phi)},\qquad
\phi_{\rm next}=\frac{P_\star}{P(\widetilde\phi)}\,\widetilde\phi.
\end{equation}

With exact group solves, the new source is proportional to \(M F\) for \(M=QA^{-1}X\).
Thus the direction of the source follows power iteration. At a dominant eigenmode
\(MF=kF\), the production ratio in \eqref{eq:iteration} returns that same \(k\). For the
connected default mesh, positive interior couplings and positive removal give
\(A_1^{-1}\) strictly positive entries. With default fast-group birth, the fast-group
production contribution \(\operatorname{diag}(\nu\Sigma_{f1})A_1^{-1}\) is strictly
positive on the fissioning cells. The thermal contribution is nonnegative because
\(A_2^{-1}\), downscatter and thermal production are nonnegative. Thus \(M\) restricted
to the fissioning cells is strictly positive. The Perron--Frobenius theorem then
supplies a simple positive dominant source mode and convergence in direction from a
positive source for the exact iteration. This argument relies on the stated connectivity
and coefficient conditions, and convergence speed depends on the dominance ratio. The
bounded, inexact implementation is assessed through its own acceptance criteria.
\cite{ref-4}

\subsection{Jacobi contraction and the acceptance
gate}\label{jacobi-contraction-and-the-acceptance-gate}

Write \(A_g=\mathcal D_g-B_g\), where \(\mathcal D_g\) is the diagonal and \(B_g\) has
nonnegative off-diagonal entries \(C_{g,ij}/V_i\). If removal is positive in every cell,
as in the default pack with the active overlays, then each diagonal exceeds the sum of
its row\textquotesingle s off-diagonal magnitudes. The Jacobi update
\(\phi^{m+1}=\phi^m+\mathcal D_g^{-1}(s-A_g\phi^m)\) has error operator
\(\mathcal D_g^{-1}B_g\) with infinity norm \cite{ref-5} ·
\href{https://github.com/hkkho/CRS/blob/4341ae8881f5827988c3d8b1859aa9e8f13cd8a6/src/ReactorSim.Core/Domain/SpatialEigenIterationContracts.cs\#L907}{Diagonal
construction}

\begin{equation}
\label{eq:jacobi}
q_g=\max_i
\frac{\sum_j C_{g,ij}}
{V_iR_{g,i}+\sum_j C_{g,ij}+\sum_b C_{g,ib}}<1 .
\end{equation}

This proves contraction for the exact-arithmetic fixed-source iteration. The equivalent
formula \(\mathcal D_g^{-1}(s+B_g\phi^m)\) also preserves nonnegative iterates for
nonnegative initial flux and source. The implementation explicitly checks finite and
nonnegative flux because it computes a residual correction in floating point. Rewriting
the inverse as a Neumann series in \(\mathcal D_g^{-1}B_g\) shows it is nonnegative, and
strictly positive between connected cells. With the fresh fuel row and no overlays, the
ratio in \eqref{eq:jacobi} is about 0.63 fast and 0.73 thermal at an interior cell with
six neighbors. These are representative row values. The outer iteration and the capped
inner solve each require a separate acceptance check.
\href{https://github.com/hkkho/CRS/blob/4341ae8881f5827988c3d8b1859aa9e8f13cd8a6/src/ReactorSim.Core/Domain/SpatialEigenIterationContracts.cs\#L681}{Jacobi
loop}

\begin{table}[!htbp]
\centering
\caption{Default convergence criteria for the embedded v11 pack.}
\label{tab:convergence}
\small
\begin{tabularx}{\linewidth}{@{}lX@{}}
\toprule
Check & Criterion \\
\midrule
Inner fixed-source residual & Absolute infinity norm \(\le 10^{-11}\) or relative
  \(\le 10^{-7}\) \\
Outer \(k\) change & Absolute \(\le 2\times10^{-7}\) or relative
  \(\le 2\times10^{-6}\) \\
Outer coupled equation residual & Relative infinity norm \(\le 2\times10^{-5}\) \\
Outer source-shape change & Infinity norm \(\le 10^{-5}\) \\
Outer relative power error & \(\lvert P-P_\star\rvert/P_\star \le 10^{-12}\) \\
\bottomrule
\end{tabularx}
\end{table}

The relative residual is \(\max\lvert L-R\rvert/\max(\lvert L\rvert+\lvert
R\rvert)\) over the relevant rows, where \(L\) and \(R\) are the
equation\textquotesingle s two sides. The outer residual recomputes both groups using
the new \(k\) and normalized flux. Source shape is the vector of cell fractions
\(V_iF_i/\sum_j V_jF_j\). The relative \(k\) change uses the larger magnitude of the old
and new eigenvalues as its scale. All four outer checks must pass together, with either
\(k\) test sufficient for that row. Caps are 128 inner iterations per group solve and
1,600 outer iterations. An exhausted solve has no usable final state. The residual
measures agreement with the assembled equations. Continuum and plant-model errors
require separate evidence.
\href{https://github.com/hkkho/CRS/blob/4341ae8881f5827988c3d8b1859aa9e8f13cd8a6/src/ReactorSim.Core/Domain/SpatialConvergenceContracts.cs\#L455}{Policy
and diagnostics} ·
\href{https://github.com/hkkho/CRS/blob/4341ae8881f5827988c3d8b1859aa9e8f13cd8a6/data/packs/candu6-two-group-diffusion-pack-v1.json}{Pack
policy}

\section{How CRS gets below five seconds}\label{how-crs-gets-below-five-seconds}

\subsection{Fewer equilibrium events}\label{fewer-equilibrium-events}

The first daily-turn implementation subdivided a simulated day into 480 intervals of 180
seconds. The archived development record reports roughly 9--11 minutes per local browser
turn for that prototype. The replacement used one frozen-flux interval and one endpoint
event. In the archived daily-mode acceptance, a two-channel day took 4.154 s and the
following empty day took 1.487 s. Those historical observations show how strongly the
workload depends on the integration cadence. Section~\ref{sec:latency}
uses the current v11 workload for the build comparison.
\href{https://github.com/hkkho/CRS/blob/2015dd334782358e910681bbd10f25d8fb3cf2d7/docs/performance/daily-turns.md}{Archived
daily-turn verification}

The public time-based mode takes one 43,200 s interval per command. When there is fuel
to move, Game stages all eight-bundle orders in canonical channel order and solves the
combined inventory in one equilibrium event. The resulting power and flux are held fixed
for a burnup update and a closed-form iodine--xenon integration in each cell. Game then
solves one endpoint equilibrium event if the run can continue. With an empty batch it
starts from the accepted shape and goes straight to exposure and the endpoint event. The
batch removes the need for one pre-exposure event per order. Freezing the shape removes
neutron-shape substeps inside the half-day and permits the analytic poison update. The
approximation leaves intrainterval flux and power unresolved. A combined post-fuel limit
violation can end the run before exposure. A numerical failure discards the detached
transaction.
\href{https://github.com/hkkho/CRS/blob/4341ae8881f5827988c3d8b1859aa9e8f13cd8a6/src/ReactorSim.Game/GameSession.cs\#L147}{Half-day
transaction} ·
\href{https://github.com/hkkho/CRS/blob/4341ae8881f5827988c3d8b1859aa9e8f13cd8a6/src/ReactorSim.Core/Domain/PracticeXenonContracts.cs\#L27}{Poison
integration} ·
\href{https://github.com/hkkho/CRS/blob/4341ae8881f5827988c3d8b1859aa9e8f13cd8a6/docs/gameplay/daily-turn-mode.md}{Shared
step semantics}

For fixed fission rate density \(f\) and flux, the poison equations in each cell have
constant coefficients. The iodine density \(N_I\) feeds xenon density \(N_{\rm Xe}\) by
decay. The two equations can be integrated over the entire interval using exponentials,
with \(\gamma\) denoting fission yields, \(\lambda\) decay constants, and
\(\sigma_{c,g}\) effective capture coefficients. The frozen flux makes the xenon removal
coefficient constant.
\href{https://github.com/hkkho/CRS/blob/4341ae8881f5827988c3d8b1859aa9e8f13cd8a6/docs/physics/xenon-calibration.md}{Poison
model and calibration}

\begin{equation}
\label{eq:poison}
\begin{aligned}
\dot N_I &= \gamma_I f-\lambda_I N_I,\\
\dot N_{\rm Xe} &= \gamma_{\rm Xe} f+\lambda_I N_I
-\big(\lambda_{\rm Xe}+\sigma_{c,1}\phi_1+\sigma_{c,2}\phi_2\big)N_{\rm Xe}.
\end{aligned}
\end{equation}

\subsection{Less work per event and per
iteration}\label{less-work-per-event-and-per-iteration}

An equilibrium event is a sequence of full-core eigensolves. The liquid-zone regulator
first solves a baseline without liquid-zone absorption and a controlled baseline with
the previous fills. It then uses a small bounded fill calculation to propose a change if
needed. Verification and one possible correction solve the actual spatial equations with
the proposed overlays. This bounds a normal event to two through four candidates. The
regulator prepares the validated inventory and base coefficient set once for the event,
then reuses those inputs and the bound edge and boundary records as absorption changes.
Each candidate still assembles its operator and diagonal from its effective
coefficients. The first candidate starts from the previous accepted or supplied flux and
\(k\), and later candidates seed from the preceding candidate. Each inner solve also
starts from the current outer flux for its group. These choices reuse a plausible shape
while checking the newly assembled equations.
\href{https://github.com/hkkho/CRS/blob/4341ae8881f5827988c3d8b1859aa9e8f13cd8a6/src/ReactorSim.Core/Domain/PracticeLiquidZoneRrsContracts.cs\#L1420}{Regulator
candidates} ·
\href{https://github.com/hkkho/CRS/blob/4341ae8881f5827988c3d8b1859aa9e8f13cd8a6/src/ReactorSim.Core/Domain/FullCoreDiffusionModelContracts.cs\#L490}{Prepared
solve} ·
\href{https://github.com/hkkho/CRS/commit/e0c839718a46dd30159c33b4785cbfd4ea421cb1}{Prepared-input
change}

The stencil stores at most six face terms per cell, and the operator compiles the
neighbor indices and conductances before iteration. A group application visits the rows
and their faces in a fixed order, with \(O(N)\) work and \(O(N)\) operator storage for
fixed stencil degree. The two group solves exploit triangular coupling and use reusable
scratch vectors. Each Jacobi loop performs one initial operator application. After a
correction, the application used to check the new residual also supplies the operator
value for the next correction. With \(m\) corrections this costs \(m+1\) applications in
place of the original \(2m\). The outer convergence check still applies both operators
to the new coupled state.
\href{https://github.com/hkkho/CRS/blob/4341ae8881f5827988c3d8b1859aa9e8f13cd8a6/src/ReactorSim.Core/Domain/SpatialOperatorContracts.cs\#L1039}{Operator
compilation} ·
\href{https://github.com/hkkho/CRS/commit/15fc41460f9920d3cf73cec399c2a4a62626ee8e}{Residual-reuse
change}

\subsection{Browser execution and the five-second
interval}\label{browser-execution-and-the-five-second-interval}

Pages builds .NET in Release with ahead-of-time compilation (AOT) to WebAssembly. The
production host runs in one Web Worker with WASM threads disabled and one CPU operator
partition. AOT changes how managed work executes. The worker keeps that work off the UI
thread. The compact response carries an authoritative snapshot patch and may also carry
a core replacement when the core changes. This reduces avoidable wire and presentation
work. \cite{ref-6} ·
\href{https://github.com/hkkho/CRS/blob/4341ae8881f5827988c3d8b1859aa9e8f13cd8a6/.github/workflows/deploy-candu-playtest.yml}{Pages
build} ·
\href{https://github.com/hkkho/CRS/blob/4341ae8881f5827988c3d8b1859aa9e8f13cd8a6/docs/spec/browser-playtest-protocol-v2.md}{Wire
protocol}

The client computes at most one endpoint ahead and blends accepted readings for
presentation over 7.5 s by default, 10 s at slow speed, or 5 s at fast speed.
Computation can overlap that display interval. If an endpoint is late, the display holds
the last solved state. The five-second interval gives an upcoming fast-playback command
a latency budget. Startup and device coverage are separate measurements.
\href{https://github.com/hkkho/CRS/blob/4341ae8881f5827988c3d8b1859aa9e8f13cd8a6/web/candu-playtest/src/DemoLoop.ts}{Playback
loop}

\section{Browser latency experiment}\label{sec:latency}

\subsection{Method}\label{method}

The experiment used the local GitHub Pages build at the \texttt{/CRS/} base path on
Windows 11 Pro (build 26200), AMD Ryzen 7 5800X3D, .NET SDK 10.0.303, and headless
Playwright Chromium 151.0.7922.34. Both builds use the same source revision, the default
v11 pack and a single-threaded Release host with profiling disabled. One build enables
AOT and the other leaves it disabled. The AOT run measures the full command stream. The
non-AOT comparison replays its first block for each seed.

For each of seeds 1001, 1002, and 1003, the benchmark initializes a new production
worker in \texttt{demo} mode. It excludes startup and one empty warm-up half-day, then
submits blocks of three compact \texttt{commit-day} commands with zero, one, and two
refuelling orders. The AOT run collects three blocks per seed and the non-AOT run
collects one. These are consecutive half-days in each run. The harness checks
acceptance, sequence, exact elapsed simulation time, and fuel consumption. It records
endpoint state digests and solver diagnostics. Samples within each seed are serially
dependent. The production AOT timings were collected before the other validation runs.
The comparison build ran while validation was in progress.
\href{https://github.com/hkkho/CRS/blob/main/web/candu-playtest/scripts/benchmark-demo-latency.mjs}{Benchmark
script}

Elapsed time starts just before serializing and posting the request and ends after
receipt and JSON parsing on the page. It includes the worker\textquotesingle s WASM
calls, any yield in the worker, snapshot mapping and serialization, transport, and
parsing. It excludes the app\textquotesingle s subsequent patch materialization,
rendering and interpolation. The worker also reports its own dispatch duration. Both
timers aggregate work across every candidate eigensolve in the command.
\href{https://github.com/hkkho/CRS/blob/4341ae8881f5827988c3d8b1859aa9e8f13cd8a6/web/candu-playtest/src/wasmWorker.ts\#L35}{Worker
timing boundary} ·
\href{https://github.com/hkkho/CRS/blob/4341ae8881f5827988c3d8b1859aa9e8f13cd8a6/web/candu-playtest/src/bridge.ts\#L451}{Application
materialization}

\subsection{Results}\label{results}

\begin{table}[!htbp]
\centering
\caption{Warm worker round-trip latency for the Release AOT build.}
\label{tab:aot}
\small
\begin{tabular*}{\linewidth}{@{\extracolsep{\fill}}lrrrr@{}}
\toprule
\shortstack[l]{Fuel orders\\per command} &
\shortstack[r]{Warm AOT\\samples} &
\shortstack[r]{AOT median\\(s)} &
\shortstack[r]{AOT maximum\\(s)} &
\shortstack[r]{AOT below\\5 s} \\
\midrule
0 & 9 & 1.913 & 2.036 & 9/9 \\
1 & 9 & 3.825 & 4.054 & 9/9 \\
2 & 9 & 3.939 & 4.201 & 9/9 \\
\bottomrule
\end{tabular*}
\end{table}

In Table~\ref{tab:aot}, one order adds the post-fuel equilibrium event to the empty-step workload. The second
order shares that event, although the perturbation and resulting iteration count can
differ. Across the AOT run, the worker timer accounted for nearly all elapsed time. The
remaining time includes message transport and parsing. Compact responses still carried
about 2.75--2.76 MB of UTF-8 JSON for these changing cores. Every measured endpoint
recorded three regulator candidate solves in its final event. The accepted endpoint
candidate took 24--70 outer iterations, and the largest reported endpoint relative
residual was \(1.9974\times10^{-5}\). Those diagnostics describe the endpoint and
exclude work on a pre-exposure event.

\begin{table}[!htbp]
\centering
\caption{Matched worker round-trip latencies for AOT and non-AOT builds.}
\label{tab:matched}
\small
\begin{tabular*}{\linewidth}{@{\extracolsep{\fill}}lrrrr@{}}
\toprule
\shortstack[l]{Fuel orders\\per command} &
\shortstack[r]{Matched samples\\per build} &
\shortstack[r]{AOT median\\(s)} &
\shortstack[r]{Non-AOT median\\(s)} &
\shortstack[r]{Non-AOT maximum\\(s)} \\
\midrule
0 & 3 & 1.689 & 12.620 & 14.421 \\
1 & 3 & 3.634 & 27.865 & 30.674 \\
2 & 3 & 4.118 & 32.244 & 34.920 \\
\bottomrule
\end{tabular*}
\end{table}

In Table~\ref{tab:matched}, the non-AOT medians were about 7.5--7.8 times the corresponding AOT medians. All
12 paired endpoints, including warm-ups, had matching state digests, endpoint candidate
counts, and endpoint outer iteration counts. This supports AOT as a major contributor to
lower execution time for this workload. The comparison uses three samples per row and
different background load, so the ratios are descriptive. The 27-command production
result supports the five-second finding on this desktop for the tested seeds and
batches. The uncompressed native WASM module was 22.6 MB with AOT and 1.50 MB without
it. Cold download and startup are separate from the warm command measurements.
\href{https://github.com/hkkho/CRS/blob/main/benchmarks/diffusion-solver-latency-2026-11-09.csv}{Per-command
samples}

An older instrumented real-time workload attributed 44.8\% of the measured
scheduled-step dispatch time to operator applications and 0.05\% to poison integration.
That profile used a different cadence and a build whose instrumentation affected
optimization. It helps explain the emphasis on reducing spatial solves and stencil
passes, while the present timings measure the production command without profiling.
\href{https://github.com/hkkho/CRS/blob/2015dd334782358e910681bbd10f25d8fb3cf2d7/docs/physics/runtime-profile.md}{Archived
runtime profile}

\subsection{Reproduction}\label{reproduction}

From the repository root, stage the runtime and run the Pages build. Start the preview
in one terminal and run the benchmark in a second from \texttt{web/candu-playtest}. To
obtain the comparison build, repeat with \texttt{-RunAOTCompilation} omitted and with
fresh MSBuild outputs for the second configuration. On Windows, a separate
\texttt{ArtifactsPath} with \texttt{UseArtifactsOutput=true} can isolate the build
intermediates. Add \texttt{1} after the URL to run one block per seed, as used for the
non-AOT comparison. The report contains raw millisecond samples and build metadata.
Isolate the timing run for a controlled comparison.
\href{https://github.com/hkkho/CRS/blob/4341ae8881f5827988c3d8b1859aa9e8f13cd8a6/docs/maintenance/hosting.md}{Build
and preview guide}

\begin{verbatim}
# From the repository root
./tools/Build-BrowserWasm.ps1 -Configuration Release `
  -RunAOTCompilation -OmitPrecompressedAssets
cd web/candu-playtest
npm run build:pages
npm run preview -- --configLoader runner --outDir dist-pages `
  --base /CRS/ --host 127.0.0.1 --port 4183 --strictPort

# From web/candu-playtest in a second terminal
node scripts/benchmark-demo-latency.mjs http://127.0.0.1:4183/CRS/
\end{verbatim}

\section{Verification and limits of
inference}\label{verification-and-limits-of-inference}

For one fresh cell with six reflective faces, leakage vanishes. With \(\chi_1=1\) and
\(\chi_2=0\), the analytic solution is

\begin{equation}
\label{eq:reflective-cell}
\begin{aligned}
\frac{\phi_2}{\phi_1} &= \frac{0.200}{0.160}=1.25,\\
k_\infty &= \frac{0.0856539290262777+1.25(0.37819461018325856)}
{0.30+0.20}=1.1167943835107017 .
\end{aligned}
\end{equation}

The regression drives the production eigen solver with the one-cell stencil and tighter
tolerances. At a 1 W target it checks \(k\) to ten decimal places, each flux to relative
error \(10^{-11}\), and relative equation and power residuals at most \(10^{-12}\). It
tests group coupling, zero leakage and normalization. Spatial tests add two- and
three-cell cases, canonical ordering, failure handling, and a 4,560-node power-sum
check. The radial and axial tests check the authored boundary conductance laws. A
full-core tolerance comparison at two fixed zone fills tests \(10^{-7}\) versus
\(10^{-8}\) inner relative tolerance. It checks reactivity differences within 0.05 mk
and cell heating differences within \(10^{-4}\) of peak reference cell power.
\href{https://github.com/hkkho/CRS/blob/4341ae8881f5827988c3d8b1859aa9e8f13cd8a6/tests/ReactorSim.Core.Tests/SingleCellReflectiveDiffusionFixtureTests.cs}{Reflective
fixture} ·
\href{https://github.com/hkkho/CRS/blob/4341ae8881f5827988c3d8b1859aa9e8f13cd8a6/tests/ReactorSim.Core.Tests/SpatialSolveTests.cs}{Spatial
tests} ·
\href{https://github.com/hkkho/CRS/blob/4341ae8881f5827988c3d8b1859aa9e8f13cd8a6/tests/ReactorSim.Core.Tests/CircularRadialBoundaryTests.cs}{Radial
tests} ·
\href{https://github.com/hkkho/CRS/blob/4341ae8881f5827988c3d8b1859aa9e8f13cd8a6/tests/ReactorSim.Core.Tests/AxialVacuumBoundaryTests.cs}{Axial
tests} ·
\href{https://github.com/hkkho/CRS/blob/4341ae8881f5827988c3d8b1859aa9e8f13cd8a6/tests/ReactorSim.Core.Tests/SpatialToleranceRegressionTests.cs}{Tolerance
regression}

Game and Browser tests check one post-batch equilibrium event, one accepted half-day
increment, deterministic replay, and rollback on numerical failure. Together these
checks cover analytic correctness for small cases, behavior of the assembled discrete
equations, and the transaction that delivers their result. The one-cell case exercises
zero leakage. The full-core tests exercise the deployed mesh, while refinement error and
independent plant accuracy remain open. The v10 geometry and v11 xenon comparisons are
calibration evidence. The frozen-source half-day approximates the time dependence
between solved endpoints.
\href{https://github.com/hkkho/CRS/blob/4341ae8881f5827988c3d8b1859aa9e8f13cd8a6/tests/ReactorSim.Game.Tests/DailyTurnTests.cs}{Game
step tests} ·
\href{https://github.com/hkkho/CRS/blob/4341ae8881f5827988c3d8b1859aa9e8f13cd8a6/tests/ReactorSim.Browser.Tests/DemoModeTests.cs}{Browser
demo tests} ·
\href{https://github.com/hkkho/CRS/blob/4341ae8881f5827988c3d8b1859aa9e8f13cd8a6/data/calibration/diffusion-v10/README.md}{Geometry
fit} ·
\href{https://github.com/hkkho/CRS/blob/4341ae8881f5827988c3d8b1859aa9e8f13cd8a6/docs/physics/xenon-calibration.md}{Xenon
fit}

The Release Core rerun passed all 131 tests, including the prepared-input comparison of
coefficient digests, eigenvalue, flux and node power. An earlier run had reported a
coefficient digest mismatch. The focused rerun and the complete rerun both passed that
assertion. On the local Pages AOT build, the direct WASM reproduction matrix matched all
six cold and warm case digests with no reported browser errors. That matrix checks
repeatability at a fixed build. The half-day benchmark uses the production worker and
tests a different command sequence. All 93 Release Game tests, the focused Browser demo
test, the Pages frontend build and all 247 Vitest tests passed. The local gameplay smoke
passed on rerun after a Mono runtime assertion during seed retry in the first attempt.
Its fuel receipt check includes queueing and playback, outside the command timer in
Section~\ref{sec:latency}.
\href{https://github.com/hkkho/CRS/blob/4341ae8881f5827988c3d8b1859aa9e8f13cd8a6/tests/ReactorSim.Core.Tests/EquilibriumCoreSolverTests.cs\#L37}{Prepared-input
regression} ·
\href{https://github.com/hkkho/CRS/blob/4341ae8881f5827988c3d8b1859aa9e8f13cd8a6/web/candu-playtest/scripts/benchmark-wasm.mjs}{Pages
reproduction harness} ·
\href{https://github.com/hkkho/CRS/blob/4341ae8881f5827988c3d8b1859aa9e8f13cd8a6/web/candu-playtest/scripts/demo-smoke.mjs}{Gameplay
smoke}

\section{Conclusion}\label{conclusion}

At the tested operating point, CRS can spend a half-day of burnup and poison exposure in
one browser command while preserving cell-resolved neutron balance at the solved
endpoints. The latency headroom between the slowest measured production command and the
fast playback interval was about 0.80 s. That headroom applies to the worker command
boundary on the measured desktop. The temporal approximation lies between equilibria,
while the coupled residual and verification cases assess the equilibria themselves.

\section*{Acknowledgment}\label{acknowledgment}
\addcontentsline{toc}{subsection}{Acknowledgment}

Gemini 4 Argon reviewed the numerical argument and the attribution of browser latency.
The review informed the distinction between exact power iteration, inexact acceptance,
and complete worker timing.

\begin{thebibliography}{9}

\bibitem{ref-1}
Chaplin, R. ``Genealogy of CANDU Reactors.'' \emph{The Essential CANDU}, Chapter 2,
UNENE, 2016.
\href{https://www.unene.ca/essentialcandu/pdf/2\%20-\%20Genealogy\%20of\%20CANDU\%20Reactors.pdf}{Chapter}

\bibitem{ref-2}
MIT OpenCourseWare. ``Criticality Conditions.'' \emph{22.05 Neutron Science and Reactor
Physics}, part 13, 2006.
\href{https://dspace.mit.edu/bitstream/handle/1721.1/74136/22-05-fall-2006/contents/lecture-notes/lecture13.pdf}{Lecture
notes}

\bibitem{ref-3}
Bernal et al. ``Resolution of the Generalized Eigenvalue Problem in the Neutron
Diffusion Equation Discretized by the Finite Volume Method.'' \emph{Abstract and Applied
Analysis}, 2014. \href{https://doi.org/10.1155/2014/913043}{doi:10.1155/2014/913043}

\bibitem{ref-4}
Boyd, S. \emph{Convex Optimization I}, lecture 6, Stanford Engineering Everywhere.
\href{https://see.stanford.edu/materials/lsocoee364a/transcripts/ConvexOptimizationI-Lecture06.html}{Perron--Frobenius
theorem}

\bibitem{ref-5}
Barrett, R. et al. \emph{Templates for the Solution of Linear Systems: Building Blocks
for Iterative Methods}. SIAM, 1994.
\href{https://www.netlib.org/templates/templates.html}{Jacobi method and convergence}

\bibitem{ref-6}
Microsoft. ``Features and configurations.'' \emph{.NET WebAssembly runtime
documentation}.
\href{https://github.com/dotnet/runtime/blob/main/src/mono/wasm/features.md\#aot}{Ahead-of-time
compilation}

\end{thebibliography}

\end{document}
